\frenchspacing
\documentclass{amsart}
\usepackage{amssymb}
\usepackage{fancyhdr}
\usepackage{bbm}
\usepackage{enumerate}
\usepackage{graphicx}
\usepackage{tikz}
\usepackage{minted}
\pagestyle{fancy}
\usepackage[margin=1in]{geometry}
\newgeometry{left=1.5cm, right=1.5cm, top = 1.5cm}
\fancyhfoffset[E,O]{0pt}
\allowdisplaybreaks


\rhead{Andrew Lys}   %% <-- your name here
\chead{Problem}
\cfoot{\thepage}
\lhead{\today}



%% your macros -->
\newcommand{\nn}{\mathbb N}    %% naturals
\newcommand{\zz}{\mathbb Z}    %%integers
\newcommand{\rr}{\mathbb R}    %% real numbers
\newcommand{\cc}{\mathbb C}    %% complex numbers
\newcommand{\ff}{\mathbb F}
\newcommand{\qq}{\mathbb Q}
\newcommand{\cA}{\mathcal{A}}
\newcommand{\cP}{\mathcal{P}}
\newcommand{\cN}{\mathcal{N}}
\newcommand{\cM}{\mathcal{M}}
\newcommand{\cB}{\mathcal{B}}
\newcommand{\cL}{\mathcal{L}}
\newcommand{\limn}{\lim_{n \to \infty}} %%lim n to infty shorthand
\newcommand{\va}{\mathbf{a}}
\newcommand{\vb}{\mathbf{b}}
\newcommand{\vc}{\mathbf{c}}
\newcommand{\vx}{\mathbf{x}}
\newcommand{\vy}{\mathbf{y}}
\newcommand{\vv}{\mathbf{v}}
\newcommand{\vw}{\mathbf{w}}
\newcommand{\vu}{\mathbf{u}}
% Source - https://tex.stackexchange.com/q/79434
% Posted by Memming, modified by community. See post 'Timeline' for change history
% Retrieved 2026-02-11, License - CC BY-SA 3.0

\newcommand\ind{\protect\mathpalette{\protect\independenT}{\perp}}
\def\independenT#1#2{\mathrel{\rlap{$#1#2$}\mkern2mu{#1#2}}}

\DeclareMathOperator{\Var}{\mathbb{Var}}  %% variance
\DeclareMathOperator{\Aut}{Aut}
\DeclareMathOperator{\Sym}{Sym}
\DeclareMathOperator{\Cov}{\mathbb{Cov}}
\DeclareMathOperator{\ee}{\mathbb{E}}
\newtheorem{theorem}{Theorem}
\theoremstyle{definition}
\newtheorem{definition}{Definition}
\newtheorem{exercise}{Exercise}


\begin{document}
\noindent
Problem Set 1  \hfill \today  %% <-- Update Notes here ***
\smallskip
\hrule
\smallskip
\noindent
Solutions by {\bf Andrew Lys} \qquad   %% <-- your name here ***
  {\tt alys(at)ncsu.edu}

\vspace{0.5cm}
\subsection*{Problem 1}
Consider 
\begin{align*}
    y &= \sigma(w^\top x)\\
    \ell &= \cL(y)
\end{align*}
for $w, x \in \rr^d$ and $\sigma(z)=z^+$. Gve an expression for $\frac{\partial \ell}{\partial w_i}$ as a function of $\frac{d \cL(y)}{dy}$.
\begin{proof}
    By the chain rule, we have:
    \begin{align*}
        \frac{\partial \ell}{\partial w_i} &= \frac{d \cL(y)}{d y}  \frac{\partial y}{\partial w_i} \\
        &= \frac{d \cL(y)}{d y}  \frac{\partial \sigma(w^\top x)}{\partial w_i} \\
        &= \frac{d \cL(y)}{d y}  \sigma'(w^\top x)  x_i\\
        &= \frac{d \cL(y)}{d y}  \mathbbm{1}_{w^\top x > 0}  x_i
    \end{align*}
\end{proof}

\subsection*{Problem 2}
\begin{align*}
    \text{for } b, j &\quad s.\text{grad}[b,j] += p.\text{grad}[b,j]  \exp(s[b,j])/Z[b]\\
    \text{for } b, j &\quad Z.\text{grad}[b] -= p.\text{grad}[b,j]  \exp(s[b,j])/(Z[b])^2\\
    \text{for } b, j & \quad s.\text{grad}[b,j] += Z.\text{grad}[b]  \exp(s[b,j])
\end{align*}

\subsection*{Problem 3}
Consider the following code:
\begin{align*}
    \text{for } b, j &\quad e[b] -= p[b,j] p.\text{grad}[b,j] \\
    \text{for } b, j &\quad s.\text{grad}[b,j] += p[b,j](p.\text{grad}[b,j] + e[b])
\end{align*}
Show how this code computes the same gradient as the code in Problem 2.
\begin{proof}
    We can plug in $Z.\text{grad}[b]$ from the second line of Problem 2 into the third line of Problem 2 to get:
    \begin{align*}
        \text{for } b, j &\quad s.\text{grad}[b,j] += \left(
            \sum_j - p.\text{grad}[b,j] \exp(s[b,j])/(Z[b])^2
        \right) \exp(s[b,j])\\
        &= \frac{e[b]}{Z[b]} \exp(s[b,j]) = p[b,j] e[b]
    \end{align*}
    Additionally, we may rewrite the first line of Problem 2 as:
    \begin{align*}
        \text{for } b, j &\quad s.\text{grad}[b,j] += p.\text{grad}[b,j]  \exp(s[b,j])/Z[b] = p[b,j] p.\text{grad}[b,j]
    \end{align*}
    Now that we have no reliance on $Z.\text{grad}[b]$, we can combine the two lines, after computing $e[b]$, to get:
    \begin{align*}
        \text{for } b, j &\quad s.\text{grad}[b,j] += p[b,j] p.\text{grad}[b,j] + p[b,j] e[b] = p[b,j](p.\text{grad}[b,j] + e[b])
    \end{align*}
\end{proof}

\subsection*{Problem 4}
\subsubsection*{Part (a)}
\begin{align*}
    \text{for } b, j, k &\quad \tilde{R}_t[b, j] += W^{h, R}[j, k] h_{t-1}[b, k]\\
    \text{for } b, j, k &\quad \tilde{R}_t[b, j] += W^{x, R}[j, k] x_t[b, k]\\
    \text{for } b, j &\quad \tilde{R}_t[b, j] -= B^R[j]
\end{align*}
\subsubsection*{Part (b)}
\begin{align*}
    \text{for } b, j &\quad B^R.\text{grad}[j] -= \frac1{B} \tilde{R}_t.\text{grad}[b, j]\\
    \text{for } b, j, k &\quad W^{x, R}.\text{grad}[j, k] += \frac1{B} \tilde{R}_t.\text{grad}[b, j] x_t[b, k]\\
    \text{for } b, j, k &\quad W^{h, R}.\text{grad}[j, k] += \frac1{B} \tilde{R}_t.\text{grad}[b, j] h_{t-1}[b, k]\\
    \text{for } b, j, k &\quad h_{t-1}.\text{grad}[b, k] +=  \tilde{R}_t.\text{grad}[b, j] W^{h, R}[j, k]\\
\end{align*}

\subsection*{Problem 5}

\begin{verbatim}
    hs = []
    hs.append(CELL(Phi, X[0], ZERO))
    for t in range(1, T + 1):
        hs.append(CELL(Phi, X[t], hs[-1]))
\end{verbatim}
\end{document}
